\documentclass[11pt]{article}
\usepackage[margin=1in]{geometry}
\usepackage{amsmath, amssymb, amsthm}
\usepackage{tcolorbox}

\newtcolorbox{definition}[1][Definition]{
    title={#1},
    colback=blue!5,
    colframe=blue!50!black,
    fonttitle=\bfseries
}

\newtcolorbox{theorem}{
    title=Theorem,
    colback=green!5,
    colframe=green!40!black,
    fonttitle=\bfseries
}

\title{Section 18 Notes}
\author{Jonathan Nguyen}
\date{}

\begin{document}
\maketitle

\begin{definition}[Ring]
A ring is a set that forms an abelian group under addition,
is closed under associative multiplication, and satisfies
the left and right distributive laws.
\end{definition}

\begin{theorem}
If $R$ is a ring, then for all $a,b \in R$ we have:
\begin{enumerate}
    \item $a0 = 0a = 0$.
    \item $a(-b) = (-a)b = -(ab)$.
    \item $(-a)(-b) = ab$.
\end{enumerate}
\end{theorem}

\begin{definition}[Homomorphism]
For ring $R$ and $R'$, a map $\phi: R \rightarrow R'$ if a homomorphism if the following two conditions are satisfied for all $a, b \in R$
\begin{enumerate}
    \item $\phi(a + b) = \phi(a) + \phi(b)$
    \item $\phi(ab) = \phi(a)\phi(b)$
\end{enumerate}
\end{definition}

\begin{definition}[Evaluation Homomorphism]
    Let $F$ be the ring of all functions $\mathbb{R} into \mathbb{R}$ defined as $(fg)(x) = f(x)g(x)$ and $(f+g)(x) = f(x) + g(x)$.
    For each $a \in \mathbb{R}$, we have the \textbf{evaluation homomorphism} $\phi_a: F \rightarrow \mathbb{R}$, where 
    $\phi(f) = f(a)$ for $f \in F$.
\end{definition}

\begin{definition}[Isomorphism]
    Isomorphism $\phi: R \rightarrow R'$ from a ring $R$ to $R'$ is a homomorphism and bijection.
\end{definition}

\begin{definition}[Commutative Ring]
    A ring is commutative if multiplcation is commutative.
\end{definition}

\begin{definition}[Ring with Unity]
    A ring with unity is a ring that contains the multiplicative identity.
\end{definition}

\begin{definition}[Unit]
    $a \in R$ is a unit if $\exists b \in R \mid ab = 1$
\end{definition}

\begin{definition}[Division Ring]
    $R$ is a division ring if every nonzero element is a unit.
\end{definition}

\begin{definition}[Field]
    A field is a commutative division ring.
\end{definition}

\begin{definition}[Subring]
    $S$ is a subring of $R$ if $S \subseteq R$ and the induced operations operate on $S$.
\end{definition}

\begin{definition}{Subfield}
    $L$ is a \textbf{subfield} if the induced operations of the field $R$ are preserved in $L$.
\end{definition}

\end{document}